% ECONOMICS_PROBLEM: {"id": "D91-638", "title": "Stability of social-preference measures", "jel_code": "D91", "source_id": "OP-0303"}
% !TeX program = xelatex
\documentclass[11pt,a4paper]{article}
\usepackage[margin=23mm]{geometry}
\usepackage{fontspec}
\setmainfont{Latin Modern Roman}
\usepackage{amsmath,amssymb,mathtools}
\usepackage{xcolor,enumitem,booktabs,longtable,array,xurl,hyperref}
\definecolor{muted}{HTML}{53636C}
\hypersetup{colorlinks=true,urlcolor=blue,linkcolor=blue}
\setlength{\parindent}{0pt}
\setlength{\parskip}{5pt}
\newcommand{\R}{\mathbb R}
\newcommand{\E}{\mathbb E}
\newcommand{\N}{\mathbb N}
\newcommand{\ind}{\mathbf 1}
\newcommand{\Var}{\operatorname{Var}}
\newcommand{\Cov}{\operatorname{Cov}}
\newcommand{\supp}{\operatorname{supp}}
\DeclareMathOperator*{\argmin}{arg\,min}
\DeclareMathOperator*{\argmax}{arg\,max}
\newcommand{\entrymeta}[3]{{\sffamily\small\color{muted}\textbf{\texttt{#1}}\quad|\quad #2\par #3\par}}
\newcommand{\Problemref}[1]{\texttt{#1}}
\newcommand{\JELref}[2]{\texttt{#2} (JEL \texttt{#1})}
\begin{document}
\section*{Stability of social-preference measures}
\textbf{Problem ID:} \texttt{D91-638}\quad \textbf{JEL:} \texttt{D91}\par
% BEGIN ECONOMICS STATEMENT
\label{problem:OP-0303}
{\sffamily\small\color{muted}Original subject group: Behavioral and experimental economics\par}
\label{definitions:problem:30007585}
\textbf{Audit classification:} Published research agenda. The 2022 manuscript and final article explicitly discuss out-of-sample validation. Final issue year is 2024, online date 2023, not 2025.
\subsection*{Definitions and precise research target}
\textbf{Model and research target.} Let \(c\) index countries or participant populations and \(k\) index trust, altruism, positive reciprocity, or negative reciprocity. For individual \(i\), observe a survey vector \(S_{ic}\), predetermined demographic covariates \(X_{ic}\), and an incentivized behavioral outcome \(Y_{ick}\) in a precisely specified game. A predictor \(f_k\), trained in a source population \(c_0\), maps survey responses and covariates to behavior. For a common loss function \(\ell_k\), define its target-population error and improvement over a covariate-only predictor \(g_k\) by
\[R_{ck}=\mathbb E_c[\ell_k(Y_{ick},f_k(S_{ic},X_{ic}))],\qquad V_{ck}=\mathbb E_c[\ell_k(Y_{ick},g_k(X_{ic}))]-R_{ck}.\]
Determine whether \(V_{ck}\) is economically meaningful across populations, and whether item rankings and conditional calibration remain stable after translation, stake normalization, and changes in social context. Estimate reliability with repeated surveys and games, separating measurement noise from population changes in conditional behavior. Assess quantitative hypothetical-game items separately from qualitative self-assessments; neither a significant source-sample correlation nor a country-average association establishes individual portability. Prespecify common coding and game protocols, quantify sampling uncertainty, and test predictors on independent target samples. The research target is a validated domain of prediction and a measurement design that improves it, rather than proof that social preferences are invariant across cultures or that any single experimental game reveals an unchanging latent trait.

\emph{Definitions and maintained assumptions (editorial unless a source definition is named).} For each country/population specify the sampling frame, translation protocol, survey coding and monetary stakes. The game outcome \(Y_{ick}\) is a bounded observable such as a transfer amount or punishment choice under a fixed game; it need not uniquely identify altruism, trust or reciprocity. Train \(f_k\) and covariate-only \(g_k\) on independent source data with the same tuning budget; use a common loss after prespecified currency/resource normalization. \(V_{ck}>0\) means predictive improvement, and a threshold \(v_0\) defines economic significance. Population calibration is, for scalar predictions, \(\mathbb E_c[Y\mid f_k(S,X)=r]=r\), interpreted almost surely; finite-sample tests use declared bins or smoothness. Repeated measurements identify test--retest stability conditional on a temporal model, not automatically a stable latent trait. Cross-country measurement invariance requires a model connecting translated items to comparable latent quantities; different survey-game correlations alone do not establish or refute it.
\subsection*{Variants and assumption changes}
\begin{enumerate}
\item \textbf{Protocol invariance (editorial).} Hold game, translated survey and normalized stakes fixed across independently sampled populations. Estimate prediction improvement and calibration with simultaneous uncertainty for all populations.
\item \textbf{Context manipulation (editorial).} Randomize partner description, anonymity or stake size within population. Determine whether survey predictions remain calibrated under these changes and distinguish context effects from measurement noise using repeated tasks.
\end{enumerate}
\subsection*{References and source audit}
\begin{enumerate}
\item IZA §1 and §4; final article §1 and §4. Supports the stated research area; exact displayed specification is editorial. \url{https://docs.iza.org/dp15006.pdf}
\item IZA §1 and §4; final article §1 and §4. Supports the stated research area; exact displayed specification is editorial. \url{https://link.springer.com/article/10.1007/s40881-023-00151-5}
\end{enumerate}
\begin{enumerate}\raggedright
\item\hypertarget{definitions:source:30007585:1}{} Michael Kosfeld; Zahra Sharafi. 2022. \emph{The Preference Survey Module: New Evidence on Social Preferences from Tehran}. IZA Discussion Paper 15006, §1, printed p. 2; §4 Conclusion, printed pp. 11–12 (PDF pp. 4 and 13–14)..
\url{https://docs.iza.org/dp15006.pdf}.
\item\hypertarget{definitions:source:30007585:2}{} Michael Kosfeld; Zahra Sharafi. 2024. \emph{The Preference Survey Module: evidence on social preferences from Tehran}. §1 Introduction, paragraphs 2–4; §4 Conclusion. Published online October 19, 2023; issue June 2024, pp. 152–164..
\url{https://link.springer.com/article/10.1007/s40881-023-00151-5}.
DOI: \url{https://doi.org/10.1007/s40881-023-00151-5}.
\end{enumerate}

\subsection*{Common conventions: Source mathematical conventions}

These conventions apply to each entry unless explicitly replaced.

\textbf{Models and identification.} A primitive model \(m\) specifies all probability laws, feasible choices, information, equilibrium and measurement rules used by its target. The admissible class \(\mathcal M\) is fixed independently of outcomes. If \(P_O(m)\) is its observable probability law and \(\tau(m)\) the target, its identified set at observable law \(P\) is
\[
 \mathcal I_\tau(P)=\{\tau(m):m\in\mathcal M,\ P_O(m)=P\}.
\]
Point identification means that this set is a singleton. An empty set signals incompatibility of data and assumptions. Coordinate bounds need not describe the joint set. Bounds are sharp only if every retained target is attainable or arbitrarily approachable by compatible models. Finite bounds require explicit restrictions on unobserved quantities; unrestricted confounding, hidden wealth, extrapolation or arbitrary equilibrium selection can yield uninformative sets.

\textbf{Causal interventions.} A potential outcome \(Y_i(p)\) is the outcome under a complete policy regime \(p\), including timing, financing, information and market spillovers. The target population and units are fixed before treatment. With interference, use the complete assignment vector or a justified exposure mapping; individual no-interference notation alone cannot identify an economy-wide intervention. A difference of mean potential outcomes does not require the cross-world joint distribution. Conditioning on post-treatment migration, survival, enrollment or employment changes the estimand and needs separate justification.

\textbf{Statistical statements.} An estimator, confidence set, forecast or test specifies a sampling law, model class and loss/coverage criterion. Uniform \((1-\alpha)\)-coverage means \(\inf_{m\in\mathcal M}\Pr_m\{\tau(m)\in C_n\}\geq1-\alpha\), or an explicitly asymptotic version. The whole parameter space trivially has coverage, so informativeness requires a stated width, risk or power objective. A forecasting improvement is relative to a named benchmark, information set and evaluation population.

\textbf{Equilibrium and optimization.} An equilibrium comprises feasible plans, optimality, beliefs where needed, budget constraints and all market/resource clearing. The primitives or a stated benchmark define each correspondence \(\mathcal E(p;m)\). Nonemptiness is assumed or proved, never inferred just from compact policy sets. A selected-equilibrium target uses a declared selection rule; a robust target quantifies over all permitted equilibria. Use suprema or infima unless attainment is established. An \(\varepsilon\)-optimal policy is feasible and within \(\varepsilon\) of the finite optimum value. Report an empty feasible set separately.

\textbf{Welfare and accounting.} Normative utility, interpersonal normalization, social weights, numeraire, discounting and the use of public receipts are primitives. Behavioral utility need not coincide with the welfare criterion. Household welfare includes ownership income and rents once; adding profits again to compensating variation generally double-counts them. A monetary equivalent is defined by an explicit transfer and price convention, with monotonicity and range conditions. Average effects, mechanism contributions, transfers and welfare losses are different targets. Infinite sums require convergence and dynamic feasible plans require terminal or no-Ponzi conditions.

\textbf{Source versions.} A working-paper date, revision date and journal publication year are distinguished. A DOI for a journal version is not automatically the DOI of an earlier manuscript. Locators refer to the stated version and printed pages where specified. A metadata-only check does not verify a claim about a full-text conclusion. Every entry's source subsection narrows the provenance claim accordingly.
% END ECONOMICS STATEMENT
\end{document}
